% !TeX root = ../教材题库_全部条目_转写工作稿.tex
% Source: XJ-B1-C04-EX-P131-Y1058
% 第4章 / 4.4 函数与方程 / 教材 p.131 / PDF 第138页
% OCR未校对：公式与文字均可能有误；未完成数学TeX转写；禁止直接发给学生
% 原题及相邻原解截图（核对用）：../题目原图/XJ-B1-C04-EX-P131-Y1058-01.webp
% 本文件只在转写完成后作为题干使用；OCR 条目尚不是可靠题干。
\begin{bankitem}{4.4 例题 例题3}{第4章 / 4.4 函数与方程 / 教材 p.131 / PDF 第138页}{XJ-B1-C04-EX-P131-Y1058}{OCR未校对：公式与文字均可能有误；未完成数学TeX转写；禁止直接发给学生}
\begin{ocrdraft}
[教材p.131 / PDF第138页]
\par
0河借助计算机讨论三次方程一32十1=0的解的个数与分布情况．
\par
解记函数不工）一於一3十1.由于不一1）一一3<0，不0）一1>0，不1）一
\par
一1<0，不2）一一3<0，不3）一1>0，可知函数丆0）在区间（一1，0），
\par
（0，1)和
\par
（2，3）内各至少有一个零点．
\par
由于丿\ensuremath{{}^\circ\mathrm C}r）一分一3十1一(x一3）十1在
\par
丿：过一丘+1
\par
（一亠一1]上为负，在[3，十。。）上为正，故丆陆）
\par
在（一亠一1]与[3，十。。）上没有零点．
\par
2
\par
借助计算机软件可作出上述函数图象，如图
\par
一2
\par
4．4一3所示．从图中可以看出，不到一分一3十1
\par
在区间（一1，0），
\par
（0，D和（2，3）内分别只有一个
\par
零点．因此，不x）在R上只有三个零点．
\par
图4．4一3
\par
综上可知，原方程有三个解，且分别位于区间
\par
（0，1）和（2，3）上．
\end{ocrdraft}
\end{bankitem}
